\subsection{Tile-Based GPU Programming}
\label{sec:bg:tile}

A GPU executes thousands of threads organized into \emph{thread blocks},
each occupying a \emph{streaming multiprocessor} (SM) with fast programmer-managed scratchpad (\emph{shared memory}).
The roofline model~\cite{williams2009roofline} captures the two performance limits:
arithmetic throughput (TFLOPS) and memory bandwidth (GB/s);
high-performance kernels tile data~\cite{wolfe1989more} to maximize shared memory reuse and stay compute-bound.

\subsection{Triton}
\label{sec:bg:triton}

Triton~\cite{tillet2019triton} is a Python-embedded DSL and compiler
in which the unit of computation is a \emph{tile}:
a statically shaped, multi-dimensional array operated on by all threads in a program instance.
The Triton compiler performs memory coalescing, shared memory allocation,
thread swizzling, and software pipelining automatically,
targeting NVIDIA PTX through an MLIR-based compilation pipeline~\cite{lattner2021mlir}.
Triton exposes performance tuning via \texttt{@triton.autotune},
sweeping programmer-specified tile configurations and caching the fastest per problem shape.
Triton has been integrated as the default lowering target for \texttt{torch.compile}
since PyTorch 2.0~\cite{li2023torchcompile}.

Triton's support for convolution is less mature:
an im2col-style lowering exists but has been observed to fall substantially short
of cuDNN performance on common workloads~\cite{triton2022conv591}.

\subsection{TileLang}
\label{sec:bg:tilelang}

TileLang~\cite{wang2025tilelang} is a DSL developed at Microsoft Research
that separates the \emph{algorithm} (what to compute) from the \emph{schedule} (how to map computation onto GPU resources) more explicitly than Triton,
allowing the programmer to independently tune memory staging, thread layout, and pipeline depth without changing the functional description.
TileLang compiles through Apache TVM's infrastructure with explicit hardware intrinsic support on NVIDIA and AMD architectures.
TileLang's convolution support and its performance relative to cuDNN
have not been systematically evaluated prior to this work.

\subsection{cuBLAS and cuDNN}
\label{sec:bg:libs}

%\paragraph{cuBLAS}
cuBLAS~\cite{cublas2024} is NVIDIA's BLAS implementation for CUDA,
providing highly optimized routines for dense linear algebra,
with \texttt{cuBLASLt} offering a flexible GEMM API
with configurable layouts and compute types (FP32, FP16, BF16, FP8, INT8).
cuBLAS dispatches the fastest GEMM implementation via a runtime heuristic recommender
trained on problem shapes and hardware generation.
The underlying template library, CUTLASS~\cite{cutlass2023},
is publicly available and serves as a reference for DSL performance comparisons.

%\paragraph{cuDNN}
cuDNN~\cite{chetlur2014cudnn,cudnn2024} covers the full set of deep learning primitives:
convolution (forward and backward passes), pooling, normalization, recurrent layers, and attention.
cuDNN selects among implicit GEMM, Winograd ($3\times3$ filters), FFT (large filters), and direct convolution at runtime via \texttt{cudnnFind}, benchmarking all candidates per configuration.

\subsection{TritonBench}
\label{sec:bg:tritonbench}

TritonBench~\cite{li2025tritonbench} is a comprehensive collection of production-grade Triton kernels
curated from high-starred open-source repositories,
covering attention variants, GEMM, softmax, normalization,
and fused operations.
It provides both functional correctness metrics and efficiency metrics
--- memory bandwidth utilization and GPU efficiency as a fraction of theoretical peak TFLOPS ---
making it a natural starting point for a performance gap study.
Our kernel suite is derived from TritonBench's
GitHub channel (184 kernels from 95 repositories),
supplemented with TileLang re-implementations and additional convolution configurations.
